由于本实验主要基于 MATLAB/Simulink 仿真，误差来源不同于实际测量实验。
主要误差来自模型参数设置、仿真步长、图形读数以及理论模型的理想化假设。

如果仿真步长过大，阶跃响应中的超调峰值和调节时间可能无法精确读出。
Bode 图和零极点图的解释也受到图像分辨率和坐标轴缩放的影响。
此外，所使用的传递函数模型忽略了实际系统中可能存在的非线性、饱和、噪声和参数不确定性。

另一类误差来自参数换算。
实验给出的 PT2 传递函数需要与标准形式 $T^2s^2+2DTs+1$ 对比，才能得到 $\omega_n$ 和 $D$。
若在换算时忽略 $T=\sqrt{a}$ 或将 $\omega_n$ 与 $T$ 混淆，就会导致对阶跃响应和极点位置的解释不一致。

仿真时间范围也会影响结论。
例如强阻尼系统响应较慢，如果仿真时间设置过短，可能误认为系统无法达到稳态。
无阻尼或弱阻尼系统则需要足够细的时间分辨率，才能准确显示振荡峰值。

因此，实验结果适合用于理解基本传递环节的典型动态特性，但不能直接等同于真实物理系统的完整行为。
尽管如此，由于本实验对象是线性传递函数模型，MATLAB 计算本身具有较高精度，主要结论不会受到这些误差来源的根本影响。
